\section{HLL and Classical Sketch Details}\label{app:min-hll}

\subsection{HLL as Mergeable State}

HyperLogLog isolates the state-versus-value distinction. The scalar answer
to a cardinality query is not mergeable: knowing \(|A|\) and \(|B|\) does
not tell us \(|A\cup B|\). The sketch is mergeable because it stores a
bounded register state carrying overlap information. Each item is hashed
to a register index and a leading-one count; the register stores the
largest count seen for that index; two summaries merge by pointwise
maximum.

On the running \emph{Hamlet} analogy, the warm-up query ``how many distinct
speaking roles appear'' is enumerable on a single play ($\approx 24$) and is
too small to motivate HLL on its own. The load-bearing query is ``how many
distinct character co-appearance pairs ever share a scene'': each scene emits
the set of unordered pairs
\(\{(c_i,c_j):i<j,\ c_i,c_j\text{ both speak}\}\); each pair is one stream
item; each act merges by taking the union of its scenes' pair-sets; the
readout is union cardinality. The merge is exactly set-union, and HLL's
pointwise-register-max approximates union cardinality without storing the
union-set itself. A single play is still the mechanism picture; the payoff
appears once the same pipeline runs on a corpus of plays, where the union of
co-appearance sets grows with the corpus while the register array stays
fixed-size.

The standard readout is
\[
\hat n=\frac{\alpha_m m^2}{\sum_j 2^{-r_j}},
\qquad
\alpha_m\approx 0.7213/(1+1.079/m),
\]
with small-range corrections. Its relative standard error is
\(1.04/\sqrt{m}\) \citep{FlajoletEtAl2007}. With \(p=14\),
\(m=16384\), so the estimator floor is about \(0.81\%\). A valid tree
compression preserves the register state and the induced readout within
that known noise floor; the floor itself is inherited from the sketch,
not introduced by the tree.

\begin{center}
\centering
\includegraphics[width=\linewidth]{hll_merge_learning_memory_median.pdf}
\captionof{figure}{\textbf{A trainable merge recovers the classical HLL
sketch at the estimator's own noise floor.} Relative RMSE against the
oracle versus HLL precision $p$ (with $m=2^p$ registers), for three
series: learned merge (median over five seeds, solid), exact HLL
(dashed), and the theoretical $1.04/\sqrt{m}$ floor (dotted). All three
curves sit within a fraction of the floor across $p\in\{6,\ldots,12\}$;
the learned merge reproduces classical behavior without a separately
tuned readout. Source:
\texttt{paper/ctreepo/scripts/regen\_paper\_hll\_figure.py}.}
\label{fig:min-hll-parity}
\end{center}

\subsection{Implementation Parity}

A native register implementation is byte-identical under flat and tree
reductions because pointwise maximum is an exact associative merge. Apache
DataSketches exposes the same cardinality semantics but may serialize an
equivalent sketch differently depending on the insertion and merge order.
The claim is therefore narrow: exact register implementations match
exactly; library sketches match at the readout level and stay inside the
estimator floor. This is the classical analogue of the C-TreePO
discipline: preserve the task state, then audit the readout.

\subsection{Learned Variants and Sketch Family}

Learned-\(g\) with a classical readout must produce register-like states
readable by the HLL estimator. Learned-\(g{+}f\) can learn a hidden state
and readout jointly. These variants are stress tests for whether a
trainable state can recover the classical sketch's behavior on its own
query; a task-specific oracle and audit is what makes a learned state
valid for a task with an unknown state.

Frequency sketches cover the remaining classical queries. Point-frequency
questions like ``how many lines does this character speak?'' or ``how often
is this theme mentioned?'' map to Count-Min
\citep{CormodeMuthukrishnan2005}. Heavy-hitter questions like ``which
characters or themes account for most of the mass?'' map to
Misra--Gries/Frequent Items
\citep{MisraGries1982,MetwallyAgrawalAbbadi2006}. These states merge by
accumulating weighted evidence and are not idempotent: repeated evidence
must count. The existing frequency benchmark serves as a parity and
provenance check for that sketch family.

\subsection{Query Taxonomy}

The useful distinction is by query. ``How many distinct speakers or
entities appear?'' is cardinality, so HLL is appropriate; so is the
sharper Hamlet variant ``how many distinct character co-appearance pairs
ever share a scene'', which is the cardinality of a union of per-scene
pair-sets. ``How many lines does this character speak?'' is point
frequency, so Count-Min is appropriate. ``Which characters account for
most lines?'' is a heavy-hitter query, so Frequent Items is appropriate.
Narrative or political importance is a task-specific oracle: the Manifesto
case runs on a learned semantic state audited by the local laws of
Section~\ref{sec:min-framework}.
